\section{The SMC Audit Chain}
\label{sec:audit}

\subsection{Output Format}

Every \smc{} run writes NDJSON (newline-delimited JSON) to the output stream.
Each line is one of three record types:

\paragraph{Particle records.}
One record per finalised particle, written in order of particle index:
\begin{verbatim}
{
  "type": "particle",
  "particle_id": <usize>,
  "stage": <u32>,
  "district_assignments": [<u8>; n],
  "log_weight": <f64>
}
\end{verbatim}

\paragraph{Resample records.}
One record per systematic resampling event:
\begin{verbatim}
{
  "type": "resample",
  "round": <u32>,
  "stage": <u32>,
  "ess_before": <f64>,
  "resample_seed": <u64>,
  "index_map": [<usize>; N]
}
\end{verbatim}
The \texttt{index\_map} field records which pre-resample particle each
post-resample particle derives from: after resampling,
$\mathbf{p}_{\rm new}[j] = \mathbf{p}_{\rm old}[\texttt{index\_map}[j]]$.
This enables complete trajectory reconstruction.

\paragraph{Metadata record.}
One record at the end of the output:
\begin{verbatim}
{
  "type": "metadata",
  "base_seed": <u64>,
  "particles": <usize>,
  "districts": <u32>,
  "tolerance": <f64>,
  "resample_threshold": <f64>,
  "resample_count": <u32>,
  "file_sha256": "<hex string>"
}
\end{verbatim}
The \texttt{file\_sha256} field is the SHA-256 hash of all preceding lines
(particle and resample records), with line endings normalised to LF.
This binds the metadata to the output file: any modification to a particle
record or resample record invalidates the hash.

\subsection{The Five-Step Verification Protocol}

A verifier who wishes to confirm that an \smc{} output file is authentic
proceeds as follows.

\begin{enumerate}
  \item \textbf{File integrity.}
        Compute the SHA-256 of all particle and resample lines (LF-normalised)
        and compare to the \texttt{file\_sha256} field in the metadata record.
        A mismatch indicates the file has been altered.

  \item \textbf{Seed spot-check.}
        For a randomly chosen particle $i$ and stage $t$, derive
        $\mathrm{particle\_seed}(s, t, i)$ from the recorded
        \texttt{base\_seed} $s$ using the SHA-256 schedule (Section~2.2),
        and confirm it matches the deterministic seed that would be used by
        \texttt{SmallRng::seed\_from\_u64} at that stage.

  \item \textbf{Resample replay.}
        For each resample record, re-derive \texttt{resample\_seed} from
        \texttt{base\_seed} and \texttt{round} using the SHA-256 schedule,
        re-run \texttt{systematic\_resample} on the pre-resample weights,
        and confirm the resulting \texttt{index\_map} matches the recorded one.

  \item \textbf{Trajectory trace.}
        For a randomly chosen particle $i$, trace its full construction
        history: starting from the post-resample lineage given by the
        \texttt{index\_map} records, identify which pre-resample particles
        contributed to particle $i$ at each stage, and confirm that the
        district assignment recorded in the particle record is consistent
        with re-running the proposal at each stage using the deterministic
        per-particle seed.

  \item \textbf{Weight validation.}
        For a randomly chosen particle $i$, recompute the log-weight by
        summing $\log|\mathcal{V}(T_i^{(t)})|$ over all stages $t$ for
        which particle $i$ was not resampled, and confirm it matches the
        recorded \texttt{log\_weight}.
\end{enumerate}

This 5-step protocol is stronger than the \flip{} audit chain~\citep{g8paper}:
a verifier can reconstruct not just the final plan, but the complete
construction history (which component was chosen, which spanning tree, which
cut edge) for any specific particle, from \texttt{base\_seed} alone plus the
\texttt{index\_map} records.

\subsection{Why the Index Map Makes Resampling Auditable}

Systematic resampling is the only non-deterministic step in the \smc{}
algorithm after conditioning on the seeds.
Without recording \texttt{index\_map}, a verifier could confirm that the
pre-resample weights were correct and that the resampling seed was derived
correctly, but could not confirm which post-resample particles correspond to
which pre-resample particles --- opening a manipulation window where an operator
swaps particles at resample time.

With \texttt{index\_map} recorded, the verifier re-runs systematic resampling
from the pre-resample weights and the recorded seed, confirms the derived
\texttt{index\_map} matches, and then traces each post-resample particle's
lineage exactly.
This closes the manipulation window.

\subsection{Integration with the SHA-256 Plan Manifest}

The \texttt{base\_seed} used for \smc{} is embedded in the broader plan
manifest SHA-256 chain~\citep{b11paper}, which links:
\begin{itemize}
  \item the \smc{} run (via \texttt{base\_seed}) to the ensemble output file;
  \item the output file (via \texttt{file\_sha256}) to the individual particle
        and resample records;
  \item the particle records (via \texttt{district\_assignments}) to the
        submitted plan.
\end{itemize}
This end-to-end chain makes it impossible to selectively manipulate any single
output particle without breaking either the \texttt{file\_sha256} binding or
the \texttt{base\_seed} derivation.
